\section{Method}\label{sec:method}

Our engine has two parts. A \emph{search} looks ahead through possible moves, and a
small \emph{neural network} judges how good a position is once the search stops
looking further. The network is trained to imitate the judgement of Scan, a strong
existing draughts engine \parencite{letouzey_scan}, which we use only as a teacher
and sparring partner; the playing engine never calls Scan. After the engine is
deployed, it keeps learning: every game it loses is turned into extra training
data. This section describes, in order, the game and its rules
(\cref{sec:rules}), how the engine stores positions and finds legal moves
(\cref{sec:board}), the network (\cref{sec:network}), the search
(\cref{sec:search}), the training data (\cref{sec:data}), how the network is
trained (\cref{sec:training}), how two versions of the engine are compared
(\cref{sec:arena}), and how the engine learns from its losses
(\cref{sec:learning}).

\subsection{The game}\label{sec:rules}

We play international draughts \parencite{fmjd_rules}. The board has $10\times10$
squares, of which only the 50 dark squares are used. They are numbered 1 to 50,
row by row, starting from the top-left dark square as seen by White. White's
twenty pieces start on squares 31--50 and Black's on squares 1--20; White moves
first. Every piece starts as a \emph{man}.

\begin{itemize}[nosep]
  \item A man moves one square diagonally forward. White men move towards
        squares 1--5, Black men towards squares 46--50.
  \item A man captures by jumping over a diagonally adjacent enemy piece onto the
        empty square directly behind it. Men may capture forwards and backwards.
  \item A \emph{king} moves any number of empty squares along a diagonal. It
        captures an enemy piece at any distance along a diagonal, provided all
        squares between them are empty, and may land on any empty square behind
        that piece.
  \item A capture may continue with further jumps by the same piece. Capturing
        is compulsory, and a player must choose a sequence that captures the
        largest possible number of pieces (men and kings count the same).
  \item Captured pieces are removed only after the whole sequence is finished,
        and no piece may be jumped twice. The capturing piece may pass through,
        or end on, its own starting square.
  \item A man becomes a king only if it \emph{ends} its move on the far row.
  \item A player with no legal move loses.
\end{itemize}

Two capture sequences with the same start square, end square and set of captured
pieces count as one move.

\paragraph{Draw rules.} We use the draw rules of the server the engine plays on,
lidraughts.org, taken from its source code \parencite{lidraughts}. The game keeps
a counter $c$ of plies (single moves by one player) since the last \emph{reset},
and a limit $L$ that depends on the material on the board:
\begin{itemize}[nosep]
  \item $L=10$ if one player has a single king and nothing else, and the other
        player has at most two pieces, at least one of them a king;
  \item $L=32$ if one player has a single king and nothing else, and the other
        player has three pieces, at least one of them a king;
  \item $L=50$ otherwise (in particular whenever more than four pieces remain).
\end{itemize}
After every move $L$ is recomputed. If $L=50$, the counter resets to 0 when the
move was a capture or moved a man (including a man that promotes), and otherwise
grows by one; so 50 plies of king moves without a capture end the game. If
$L<50$, the counter resets only for the move that \emph{created} the endgame
and otherwise grows by one, whatever is moved. A move creates the endgame if it
is a capture made while more than 3 pieces (for $L=10$) or more than 4 pieces
(for $L=32$) were on the board, or a promotion after which the moving player has
exactly one king. The game
is drawn when $c \ge L$, or when the same position, with the same player to move,
has occurred three times since the last reset.

\subsection{Positions and move generation}\label{sec:board}

A position is stored as three 64-bit integers: one bit set for every white piece,
one for every black piece, and one for every king, plus the side to move. Square
$s$ is stored in bit
\begin{equation}\label{eq:bit}
  b(s) = (s-1) + \left\lfloor \tfrac{1}{2}\left\lfloor \tfrac{s-1}{5}\right\rfloor \right\rfloor ,
\end{equation}
which leaves one unused bit after every two rows (bits 10, 21, 32 and 43). With
this layout the four diagonal neighbours of any square are always at bit offsets
$-6$, $-5$, $+5$ and $+6$, and a step that would leave the board lands on an
unused bit or outside bits 0--53. Each position also carries a 64-bit hash key,
the exclusive-or of a fixed random number for every (colour, piece type, square)
that is occupied and one for the side to move \parencite{zobrist1970}.

Legal moves are generated as follows. First, for every own piece, all capture
sequences are found by depth-first search: from the current square, try each of
the four directions; for a man, the neighbouring square must hold an enemy piece
not yet captured in this sequence and the square behind it must be empty; for a
king, slide over empty squares to the first piece, which must be such an enemy
piece, and then try every empty landing square behind it. During the search the
moving piece is treated as absent from its start square, and already-captured
pieces stay on the board as obstacles. A sequence ends when no further jump is
possible. Only sequences with the largest number of captures are kept, and
duplicates (same start, end and captured set) are removed. If no capture exists,
the legal moves are the ordinary moves of men and kings.

The move generator was checked in two ways: the number of move sequences from
the start position (``perft'') matches published values up to depth 9
(41{,}022{,}423 sequences), and its legal moves agree with an independent
implementation, the pydraughts library, on every position of many random games.

\subsection{The evaluation network}\label{sec:network}

\paragraph{Input.} The network looks at the board twice: once from White's point
of view and once from Black's. For a point of view $p$ we build a vector
$x_p \in \{0,1\}^{200}$ with one entry for every combination of piece type
$t$ and square:
\[
  t = 0:\ \text{own man},\quad t = 1:\ \text{own king},\quad
  t = 2:\ \text{opponent man},\quad t = 3:\ \text{opponent king},
\]
where ``own'' means belonging to player $p$. Black's view is the board turned
by 180 degrees, so that both players' men always move towards the same end of
their own view. A piece on square $s$ therefore sets entry
\begin{equation}\label{eq:feature}
  i = 50\,t + r - 1, \qquad r = \begin{cases} s & \text{for White's view},\\ 51 - s & \text{for Black's view},\end{cases}
\end{equation}
of $x_p$ to 1. All other entries are 0.

\paragraph{Layers.} Write $\clip(z) = \min(\max(z,0),1)$, applied to every entry
of a vector, and $\sigm(z) = 1/(1+e^{-z})$. Let $m$ be the player to move and
$n$ the other player. The network computes
\begin{align}
  a_p &= W_1 x_p + b_1 \quad \text{for } p \in \{m, n\}, \label{eq:l1}\\
  h_1 &= \clip\bigl([\,a_m;\ a_n\,]\bigr), \label{eq:concat}\\
  h_2 &= \clip(W_2 h_1 + b_2), \label{eq:l2}\\
  h_3 &= \clip(W_3 h_2 + b_3), \label{eq:l3}\\
  o   &= w_4^{\top} h_3 + b_4 . \label{eq:out}
\end{align}
The first layer $W_1 \in \mathbb{R}^{256\times200}$ is shared by both points of
view, so $a_m$ and $a_n$ are 256 numbers each. In \cref{eq:concat} the two are
stacked, the player to move first, giving 512 inputs to the second layer. The
remaining sizes are $W_2 \in \mathbb{R}^{32\times512}$,
$W_3 \in \mathbb{R}^{32\times32}$ and $w_4 \in \mathbb{R}^{32}$.

\paragraph{Output.} The output $o$ is read as the chance that the player to move
wins, $\sigm(o)$. Inside the search the engine needs a score in hundredths of a
man (``centi-men''); it uses
\begin{equation}\label{eq:score}
  v = \operatorname{trunc}(K\,o), \quad \text{limited to } [-5000,\,5000],
\end{equation}
where $K$ is the scale fixed during training (\cref{eq:k}). With $K=130$, a score
of $+130$ corresponds to $\sigm(1)\approx 0.73$: the player to move is expected to
score about 73\% (a draw counting as half a win).

\paragraph{Fast evaluation during search.} The vectors $a_m, a_n$ of
\cref{eq:l1} are the expensive part, but a move changes only a few entries of
$x_p$. The engine therefore keeps $a_{\text{White}}$ and $a_{\text{Black}}$ for
every position on the search path and updates them when a move is made: for each
point of view it subtracts the column of $W_1$ belonging to the moving piece on
its start square, adds the column for the piece on its end square (as a king if
it promotes), and subtracts the column of every captured piece. The result is
identical to computing \cref{eq:l1} from scratch. The second layer
(\cref{eq:l2}) is computed with 8-bit integers on processors that support it:
each input is rounded to $q = \operatorname{round}(255\,h_1) \in \{0,\dots,255\}$,
each weight to $Q = \operatorname{round}(s\,W_2)$ with $s = 127/\max|W_2|$, and the
pre-activation is recovered as $\sum q\,Q / (255\,s) + b_2$. All other steps use
32-bit floating point. A test checks that the engine's scores differ from those
of the training code by at most 3 centi-men or 2\% of the largest score in the
test, whichever is larger.

\subsection{Search}\label{sec:search}

To choose a move the engine runs a negamax alpha-beta search
\parencite{knuth1975}: it explores the tree of possible continuations, always
from the point of view of the player to move, and skips branches that cannot
change the result. The engine receives the game as the start position plus the
list of all moves played, and replays them to know the draw counter and which
positions have occurred since the last reset. Scores are in centi-men; a draw scores 0, and a position where
the player to move has no legal move at $d$ plies from the root scores
$-(10000 - d)$, so that faster wins score higher.

\paragraph{Iterative deepening.} The search is run to depth 1, then 2, 3 and so
on, each time reusing what earlier iterations stored in a hash table. From depth
4 on, each iteration first searches a narrow window of $\pm30$ around the previous
iteration's score; if the true score falls outside, the window is widened on that
side by doubling its width until it fits. If the root has only one legal move,
the engine plays it at once.

\paragraph{One node of the search.} At a node at distance $d$ from the root, with
remaining depth $D$ and score window $(\alpha,\beta)$, the engine does the
following:
\begin{enumerate}[nosep]
  \item If the position is drawn by the rules of \cref{sec:rules}, return 0. In
        the search any repetition of a position since the last reset already
        counts as a draw.
  \item If $D \le 0$, return the quiescence search below.
  \item Narrow $(\alpha,\beta)$ to the scores still possible at distance $d$
        (no win can come sooner than at $d+1$ and no loss sooner than at $d$);
        if the window becomes empty, return $\alpha$.
  \item Look the position up in the hash table. Outside the main line, if the
        stored result came from a search at least $D$ deep and settles the
        window, return it.
  \item Generate the legal moves. With none, return $-(10000-d)$.
  \item If there is exactly one legal move, search it one ply deeper
        ($D$ is not reduced).
  \item Order the moves: the move stored in the hash table first; then
        captures (those taking more kings first); then two ``killer'' moves that
        recently caused a cut-off at this distance from the root
        \parencite{akl1977}; then the rest by a history score
        \parencite{schaeffer1989}.
  \item Search the first move with the full window and every later move with a
        zero-width window around $\alpha$, re-searching with the full window if
        it turns out better \parencite{marsland1982}. A non-capturing,
        non-promoting later move is searched 1 ply shallower when $D\ge2$ and
        it is at least the 2nd move in the order (the 4th on the main line),
        and 2 plies shallower when it is at least the 5th move off the main
        line; if such a move scores above $\alpha$, it is searched again at
        full depth.
  \item As soon as a move scores at least $\beta$, stop. If it was not a
        capture, make it this node's first killer move (the old first killer
        becomes the second) and add $D^2$ to its history score.
  \item Store the best score, the best move, $D$ and whether the score is exact,
        a lower bound or an upper bound in the hash table.
\end{enumerate}
The hash table is organised in groups of five entries; a new entry replaces the
entry for the same position, else an empty one, else the one with the smallest
depth minus four times its age (in searches). History scores are divided by 8
at the start of every search and halved when one exceeds $2^{20}$.

\paragraph{Quiescence search.} At depth 0 the position may still be in the middle
of an exchange, so the network is not asked yet. If the player to move must
capture, every capture is searched (there is no option to stop, because
capturing is compulsory). Otherwise the network score $e$ is taken. It is
returned as the result if $e \ge \beta$, if the opponent could not capture
anything were it their turn, or if this is not the first quiet position reached
in the quiescence search. In the remaining case (the opponent threatens a capture
right at the horizon) every legal move is searched one ply further, with $e$ as
the score to beat.

\paragraph{Time.} Given the remaining clock time $T$ (in seconds) and increment
$I$, the engine plans to use $t_{\text{soft}} = \min\bigl((T-0.15)/30 + 0.75\,I,\
(T-0.15)/4\bigr)$ and never more than $t_{\text{hard}} =
\min(3.5\,t_{\text{soft}},\ 0.4\,(T-0.15))$. A new iteration is not started
after $0.6\,t_{\text{soft}}$, and a running one is abandoned at
$t_{\text{hard}}$, in which case the result of the last finished iteration is
played.

\paragraph{Live play.} On lidraughts the engine runs with four threads. Three
helper threads search the same position without a time limit and share the hash
table; the move comes from the main thread, and the helpers stop when it
finishes. The engine also thinks during the opponent's turn on the move it
expects. To avoid playing the same opening in every game, during the first 10
plies of a game, and only if less than 80\% of the planned time
$t_{\text{soft}}$ has been used, the engine re-scores every legal move with a search of depth
$\min(D_{\text{last}}-2,\ 8)$, where $D_{\text{last}}$ is the depth of its last
finished iteration, and plays a uniformly random move among those within 15
centi-men of the best.

\subsection{Training data}\label{sec:data}

Every training example is a position together with two labels, both from the
point of view of the player to move: a \emph{teacher score} $S$ in centi-men
(Scan's evaluation after searching the position, limited to $[-3000,3000]$) and
the \emph{game result} $R\in\{+1,0,-1\}$ (win, draw or loss of the game the
position came from). We only keep \emph{quiet} positions, in which the player to
move has no capture and more than one legal move; positions in the middle of an
exchange are left to the search.

\paragraph{Teacher self-play.} Scan plays games against itself, with its opening
book switched off (as in all our uses of Scan). Each game starts
with $k$ uniformly random legal moves, with $k$ drawn uniformly from
$\{2,\dots,8\}$. Then, on every move, Scan searches the position for a fixed
number of nodes $N$; the position is recorded with Scan's score if it is quiet,
and Scan's move is played, except that with probability 0.01 a uniformly random
legal move is played instead. A game ends when a player has no legal move, by
the draw rules of \cref{sec:rules}, as a draw after 300 plies, or as a win when
Scan's score has been at least 1500 in favour of the same player for 6
consecutive plies. We generated two sets: 400{,}000 games with $N=100{,}000$
(44.6 million positions) and a second set with $N=1{,}000{,}000$, a ten times
deeper search (7.3 million positions when the current network was trained).

\paragraph{Human games.} We downloaded rated games of the 100 highest-rated
players on each lidraughts leaderboard (blitz, bullet, rapid, classical), up to
their 1000 most recent games each, keeping games of standard draughts from the
normal start position. From games in which both players were rated at least 1800
we took every quiet position, scored it with Scan at $N=100{,}000$, and attached
the game's result (a loss on time counts as a loss), which gave about 1.3
million positions.

\paragraph{Lessons.} Games the engine lost produce further examples, described
in \cref{sec:learning}.

\subsection{Training}\label{sec:training}

\paragraph{Score scale.} The teacher score $S$ is turned into a win probability
by $\sigm(S/K)$. The scale $K$ is chosen once, as the value in
$\{50, 55, \dots, 795\}$ for which $\sigm(S/K)$ best predicts the game results
of the first two million training positions:
\begin{equation}\label{eq:k}
  K = \arg\min_{K} \; -\frac{1}{N}\sum_{j=1}^{N}
      \Bigl[ y_j \log \sigm(S_j/K) + (1-y_j) \log\bigl(1-\sigm(S_j/K)\bigr) \Bigr],
  \qquad y_j = \tfrac{R_j+1}{2}.
\end{equation}
This gives $K=130$.

\paragraph{Target and loss.} Each position gets the target
\begin{equation}\label{eq:target}
  y^{*} = \lambda\, \sigm(S/K) + (1-\lambda)\, \frac{R+1}{2}, \qquad \lambda = 0.75,
\end{equation}
a mix of mostly the teacher's opinion and partly the actual outcome. The network
is trained to make $\sigm(o)$ (\cref{eq:out}) close to $y^{*}$, using the mean
squared error $\bigl(\sigm(o)-y^{*}\bigr)^2$ averaged over a batch.

\paragraph{Procedure.} Training uses the Adam optimiser \parencite{kingma2015}
with batches of 16{,}384 positions for 15 passes over the data. The learning rate
follows a one-cycle schedule \parencite{smith2019}: it rises to $10^{-3}$ over the
first 5\% of steps and then decreases. The data are stored in files of whole
games; 2\% of the files are held out for validation, so no game is split between
training and validation, and after every pass the network with the lowest
validation loss so far is kept. A data source can be given a repeat factor $k$,
which makes each of its positions appear $k$ times per pass. The network
currently playing was trained on the 44.6 million teacher positions (factor 1)
together with the human positions, the deeper teacher positions and the lesson
positions (factor 2 each).

\subsection{Comparing two engines}\label{sec:arena}

Two versions of the engine are compared by playing matches. Each opening is a
sequence of 4 uniformly random legal moves from the start position (each opening
position distinct), and each opening is played twice with colours swapped.
Both sides play on a clock of 5 seconds plus 0.05 seconds per move. A game is
lost by a player with no legal move or a player whose clock falls below
$-0.5$ seconds, and is drawn by the draw rules or after 300 plies. The match score
$P$ is the average over games of 1 for a win, $\tfrac12$ for a draw and 0 for a
loss, and the strength difference in Elo is
\begin{equation}\label{eq:elo}
  \Delta = -400 \log_{10}\!\left(\frac{1}{P} - 1\right).
\end{equation}

\subsection{Learning from losses}\label{sec:learning}

\paragraph{Lessons from a lost game.} Whenever the engine loses a game against a
human on lidraughts, the following runs automatically.
\begin{enumerate}[nosep]
  \item \emph{Analysis.} Scan scores every position of the game, searching
        2{,}000{,}000 nodes each.
  \item \emph{Mistakes.} One of the engine's moves is a mistake if Scan would
        have played a different move and Scan's score drops by at least 0.3 men
        across it: the score before the move minus the (sign-flipped) score of
        the position after it.
  \item \emph{Lesson data.} The lesson contains every quiet position of the game
        with its 2{,}000{,}000-node score and the game's result, plus Scan
        self-play games (as in \cref{sec:data}, with $N=100{,}000$, no random
        opening moves and probability 0.03 of a random move) started 13 times
        from each mistake position and 13 times from each of the two positions
        before it.
  \item \emph{Fine-tuning.} Starting from the current network, train for 3
        passes with a peak learning rate of $2\times10^{-4}$, the same $K$ and
        $\lambda$, on 200 randomly chosen files of teacher self-play data plus
        the lesson with repeat factor 8.
  \item \emph{Check.} The fine-tuned network plays a 200-game match against the
        current one (\cref{sec:arena}). It replaces the current network only if
        it scores at least 49\%, so that a single game cannot make the engine
        noticeably weaker. The engine loads the network afresh at the start of
        each game.
\end{enumerate}

\paragraph{Sparring against the teacher.} The same procedure is also run in
rounds against Scan. In each round the current network plays 100 games against
Scan (\cref{sec:arena}). The first 12 lost games are analysed as above,
ignoring the 4 random opening moves; games without a mistake give no lesson. The network is then fine-tuned on all
lessons collected so far, from sparring and from human games, each with repeat
factor 4, plus 600 randomly chosen files of teacher self-play data. The result
replaces the current network only if it scores at least 52.5\% in a 400-game
match against it. Each network that becomes the current one is measured once
against Scan on the same 400 games (fixed openings), so that networks can be
compared with each other.
